% Options for packages loaded elsewhere
\PassOptionsToPackage{unicode}{hyperref}
\PassOptionsToPackage{hyphens}{url}
\documentclass[
  11pt,
  a4paper,
]{article}
\usepackage{xcolor}
\usepackage[margin=18mm]{geometry}
\usepackage{amsmath,amssymb}
\setcounter{secnumdepth}{-\maxdimen} % remove section numbering
\usepackage{iftex}
\ifPDFTeX
  \usepackage[T1]{fontenc}
  \usepackage[utf8]{inputenc}
  \usepackage{textcomp} % provide euro and other symbols
\else % if luatex or xetex
  \usepackage{unicode-math} % this also loads fontspec
  \defaultfontfeatures{Scale=MatchLowercase}
  \defaultfontfeatures[\rmfamily]{Ligatures=TeX,Scale=1}
\fi
\ifPDFTeX\else
  % xetex/luatex font selection
    \setmainfont[]{Noto Serif}
    \setsansfont[]{Noto Sans}
    \setmonofont[]{Noto Sans Mono}
  \setmathfont[]{Noto Sans Math}
\fi
% Use upquote if available, for straight quotes in verbatim environments
\IfFileExists{upquote.sty}{\usepackage{upquote}}{}
\IfFileExists{microtype.sty}{% use microtype if available
  \usepackage[]{microtype}
  \UseMicrotypeSet[protrusion]{basicmath} % disable protrusion for tt fonts
}{}
\makeatletter
\@ifundefined{KOMAClassName}{% if non-KOMA class
  \IfFileExists{parskip.sty}{%
    \usepackage{parskip}
  }{% else
    \setlength{\parindent}{0pt}
    \setlength{\parskip}{6pt plus 2pt minus 1pt}}
}{% if KOMA class
  \KOMAoptions{parskip=half}}
\makeatother
\ifLuaTeX
\usepackage[bidi=basic]{babel}
\else
\usepackage[bidi=default]{babel}
\fi
\babelprovide[main,import]{russian}
\ifPDFTeX
\else
\babelfont{rm}[]{Noto Serif}
\fi
% get rid of language-specific shorthands (see #6817):
\let\LanguageShortHands\languageshorthands
\def\languageshorthands#1{}
\setlength{\emergencystretch}{3em} % prevent overfull lines
\providecommand{\tightlist}{%
  \setlength{\itemsep}{0pt}\setlength{\parskip}{0pt}}
\usepackage{bookmark}
\IfFileExists{xurl.sty}{\usepackage{xurl}}{} % add URL line breaks if available
\urlstyle{same}
\hypersetup{
  pdftitle={Билеты по дискретным случайным процессам},
  pdflang={ru-RU},
  hidelinks,
  pdfcreator={LaTeX via pandoc}}

\title{Билеты по дискретным случайным процессам}
\author{}
\date{}

\begin{document}
\maketitle

\section{Билет 01. Системы множеств: алгебры, полуалгебры,
сигма-алгебры. Меры Дирака и
Жордана}\label{ux431ux438ux43bux435ux442-01.-ux441ux438ux441ux442ux435ux43cux44b-ux43cux43dux43eux436ux435ux441ux442ux432-ux430ux43bux433ux435ux431ux440ux44b-ux43fux43eux43bux443ux430ux43bux433ux435ux431ux440ux44b-ux441ux438ux433ux43cux430-ux430ux43bux433ux435ux431ux440ux44b.-ux43cux435ux440ux44b-ux434ux438ux440ux430ux43aux430-ux438-ux436ux43eux440ux434ux430ux43dux430}

Источники: Ширяев 1, гл. II; лекции ДСП.

\subsection{1. Короткий
ответ}\label{ux43aux43eux440ux43eux442ux43aux438ux439-ux43eux442ux432ux435ux442}

В теории вероятностей в общем случае нельзя требовать, чтобы хорошая
вероятность была задана на всех подмножествах пространства исходов.
Сначала выбирают систему множеств, устойчивую относительно нужных
операций: алгебру, полуалгебру или \(\sigma\)-алгебру. Алгебра замкнута
относительно конечных операций, \(\sigma\)-алгебра - относительно
счетных операций, а полуалгебра удобна как простая область задания меры
перед ее продолжением. Два базовых примера мер: мера Дирака

\[
\delta_x(A)=\mathbf 1_A(x)
\]

и мера Жордана, которая на прямоугольнике \(I=\prod_{k=1}^d[a_k,b_k)\)
задается формулой

\[
v(I)=\prod_{k=1}^d(b_k-a_k).
\]

\subsection{2. Полный
ответ}\label{ux43fux43eux43bux43dux44bux439-ux43eux442ux432ux435ux442}

\subsubsection{2.1. Пространство элементарных исходов и
события}\label{ux43fux440ux43eux441ux442ux440ux430ux43dux441ux442ux432ux43e-ux44dux43bux435ux43cux435ux43dux442ux430ux440ux43dux44bux445-ux438ux441ux445ux43eux434ux43eux432-ux438-ux441ux43eux431ux44bux442ux438ux44f}

Пусть фиксировано множество \(\Omega\). Его элементы называются
элементарными исходами, а подмножества \(\Omega\) интерпретируются как
события. Однако в общем бесконечном случае нельзя заранее рассчитывать,
что нужная вероятность будет корректно задана на всех подмножествах.
Поэтому выбирают класс допустимых множеств
\(\mathcal F\subset 2^\Omega\) и требуют, чтобы он был устойчив
относительно естественных операций над событиями.

\subsubsection{2.2. Требования к системам
множеств}\label{ux442ux440ux435ux431ux43eux432ux430ux43dux438ux44f-ux43a-ux441ux438ux441ux442ux435ux43cux430ux43c-ux43cux43dux43eux436ux435ux441ux442ux432}

Минимальные требования зависят от того, какие операции мы хотим
выполнять. Если нужно брать только конечные объединения, пересечения и
дополнения, достаточно алгебры. Если нужно работать с пределами
последовательностей событий, например с событиями вида
\(\bigcup_{n=1}^{\infty}A_n\), \(\bigcap_{n=1}^{\infty}A_n\),
\(\limsup A_n\), нужна \(\sigma\)-алгебра. Полуалгебра обычно
используется как ``кирпичики'': на ней легко задать меру, например длину
полуинтервала, а затем продолжить меру на алгебру и дальше на
\(\sigma\)-алгебру.

\subsubsection{2.3. Полуалгебры и алгебры элементарных
множеств}\label{ux43fux43eux43bux443ux430ux43bux433ux435ux431ux440ux44b-ux438-ux430ux43bux433ux435ux431ux440ux44b-ux44dux43bux435ux43cux435ux43dux442ux430ux440ux43dux44bux445-ux43cux43dux43eux436ux435ux441ux442ux432}

Для вероятностей главная система множеств - \(\sigma\)-алгебра событий
\(\mathcal F\). Но перед ней часто стоят более простые классы. Например,
на \(\mathbb R\) естественно начать с полуинтервалов \([a,b)\), на
\(\mathbb R^d\) - с полуоткрытых прямоугольников
\(\prod_{k=1}^d[a_k,b_k)\). Они образуют полуалгебру, потому что
пересечение двух таких прямоугольников снова прямоугольник такого же
вида, а разность прямоугольников раскладывается в конечное дизъюнктное
объединение прямоугольников. Из полуалгебры берут конечные дизъюнктные
объединения и получают алгебру элементарных множеств.

\subsubsection{2.4. Понятие меры и свойство
аддитивности}\label{ux43fux43eux43dux44fux442ux438ux435-ux43cux435ux440ux44b-ux438-ux441ux432ux43eux439ux441ux442ux432ux43e-ux430ux434ux434ux438ux442ux438ux432ux43dux43eux441ux442ux438}

Мера на такой системе множеств - это функция, которая каждому
допустимому множеству ставит в соответствие его ``размер''. В этом
билете достаточно понимать меру как неотрицательную конечно-аддитивную
функцию на выбранной системе множеств. Счетная аддитивность будет
центральной темой следующего билета. Для попарно непересекающихся
допустимых множеств \(A_1,\ldots,A_n\) конечная аддитивность означает

\[
\mu\left(\bigsqcup_{k=1}^n A_k\right)=\sum_{k=1}^n\mu(A_k).
\]

\subsubsection{2.5. Мера
Дирака}\label{ux43cux435ux440ux430-ux434ux438ux440ux430ux43aux430}

Мера Дирака \(\delta_x\) сосредоточивает всю массу в одной точке \(x\).
Она равна \(1\), если множество содержит \(x\), и равна \(0\), если не
содержит. Это простейшая вероятностная мера: вся случайность вырождается
в один исход.

\subsubsection{2.6. Мера Жордана как геометрический
объем}\label{ux43cux435ux440ux430-ux436ux43eux440ux434ux430ux43dux430-ux43aux430ux43a-ux433ux435ux43eux43cux435ux442ux440ux438ux447ux435ux441ux43aux438ux439-ux43eux431ux44aux435ux43c}

Мера Жордана - геометрический объем элементарных множеств. На прямой это
длина, на плоскости площадь, в \(\mathbb R^d\) - \(d\)-мерный объем.
Сначала ее задают на прямоугольниках формулой \(v(I)=\prod(b_k-a_k)\),
затем на элементарном множестве \(A=\bigsqcup_{j=1}^m I_j\) кладут

\[
v(A)=\sum_{j=1}^m v(I_j).
\]

\subsubsection{2.7. Корректность меры
Жордана}\label{ux43aux43eux440ux440ux435ux43aux442ux43dux43eux441ux442ux44c-ux43cux435ux440ux44b-ux436ux43eux440ux434ux430ux43dux430}

Основной вопрос здесь: не зависит ли эта сумма от способа разбиения
\(A\) на прямоугольники? Ответ: не зависит. Это доказывается переходом к
общему измельчению двух разбиений. Поэтому мера Жордана корректно
определена на алгебре элементарных множеств.

\subsection{3. Основные
определения}\label{ux43eux441ux43dux43eux432ux43dux44bux435-ux43eux43fux440ux435ux434ux435ux43bux435ux43dux438ux44f}

\subsubsection{Система
множеств}\label{ux441ux438ux441ux442ux435ux43cux430-ux43cux43dux43eux436ux435ux441ux442ux432}

Любое семейство подмножеств \(\mathcal C\subset 2^\Omega\).

\subsubsection{Алгебра
множеств}\label{ux430ux43bux433ux435ux431ux440ux430-ux43cux43dux43eux436ux435ux441ux442ux432}

Семейство \(\mathcal A\subset 2^\Omega\) называется алгеброй на
\(\Omega\), если:

\begin{enumerate}
\def\labelenumi{\arabic{enumi}.}
\tightlist
\item
  \(\Omega\in\mathcal A\);
\item
  если \(A\in\mathcal A\), то
  \(\overline{A}=\Omega\setminus A\in\mathcal A\);
\item
  если \(A,B\in\mathcal A\), то \(A\cup B\in\mathcal A\).
\end{enumerate}

Из этих условий сразу следует, что \(\varnothing\in\mathcal A\),
\(A\cap B\in\mathcal A\), \(A\setminus B\in\mathcal A\), а также что
алгебра замкнута относительно любых конечных объединений и конечных
пересечений.

\subsubsection{\texorpdfstring{\(\sigma\)-алгебра}{\textbackslash sigma-алгебра}}\label{sigma-ux430ux43bux433ux435ux431ux440ux430}

Семейство \(\mathcal F\subset 2^\Omega\) называется \(\sigma\)-алгеброй,
если:

\begin{enumerate}
\def\labelenumi{\arabic{enumi}.}
\tightlist
\item
  \(\Omega\in\mathcal F\);
\item
  \(A\in\mathcal F\Rightarrow \overline{A}\in\mathcal F\);
\item
  для любой последовательности \(A_1,A_2,\ldots\in\mathcal F\) выполнено
  \(\bigcup_{n=1}^{\infty}A_n\in\mathcal F\).
\end{enumerate}

Любая \(\sigma\)-алгебра является алгеброй, но не всякая алгебра
является \(\sigma\)-алгеброй.

\subsubsection{Полуалгебра}\label{ux43fux43eux43bux443ux430ux43bux433ux435ux431ux440ux430}

Семейство \(\mathcal S\subset 2^\Omega\) называется полуалгеброй, если:

\begin{enumerate}
\def\labelenumi{\arabic{enumi}.}
\tightlist
\item
  \(\varnothing\in\mathcal S\);
\item
  если \(A,B\in\mathcal S\), то \(A\cap B\in\mathcal S\);
\item
  если \(A,B\in\mathcal S\), то разность \(A\setminus B\) представима
  как конечное дизъюнктное объединение множеств из \(\mathcal S\):
\end{enumerate}

\[
A\setminus B=\bigsqcup_{k=1}^n C_k,\qquad C_k\in\mathcal S.
\]

Иногда дополнительно требуют \(\Omega\in\mathcal S\); в вероятностных
приложениях это удобно, но главное свойство полуалгебры - хорошее
поведение пересечений и разностей.

\subsubsection{Мера на
алгебре}\label{ux43cux435ux440ux430-ux43dux430-ux430ux43bux433ux435ux431ux440ux435}

Неотрицательная функция \(\mu:\mathcal A\to[0,\infty]\) называется
конечно-аддитивной мерой, если \(\mu(\varnothing)=0\) и для любых
попарно непересекающихся \(A_1,\ldots,A_n\in\mathcal A\), таких что
\(\bigsqcup_{k=1}^nA_k\in\mathcal A\), выполнено

\[
\mu\left(\bigsqcup_{k=1}^nA_k\right)=\sum_{k=1}^n\mu(A_k).
\]

\subsubsection{Мера
Дирака}\label{ux43cux435ux440ux430-ux434ux438ux440ux430ux43aux430-1}

Пусть \(x\in\Omega\). На любой \(\sigma\)-алгебре
\(\mathcal F\subset 2^\Omega\) мера Дирака в точке \(x\) задается
формулой

\[
\delta_x(A)=
\begin{cases}
1, & x\in A,\\
0, & x\notin A.
\end{cases}
\]

\subsubsection{\texorpdfstring{Полуоткрытый прямоугольник в
\(\mathbb R^d\)}{Полуоткрытый прямоугольник в \textbackslash mathbb R\^{}d}}\label{ux43fux43eux43bux443ux43eux442ux43aux440ux44bux442ux44bux439-ux43fux440ux44fux43cux43eux443ux433ux43eux43bux44cux43dux438ux43a-ux432-mathbb-rd}

Множество вида

\[
I=\prod_{k=1}^d[a_k,b_k),\qquad a_k\le b_k.
\]

Если хотя бы для одного \(k\) выполнено \(a_k=b_k\), прямоугольник
вырожден и имеет объем \(0\).

\subsubsection{Элементарное
множество}\label{ux44dux43bux435ux43cux435ux43dux442ux430ux440ux43dux43eux435-ux43cux43dux43eux436ux435ux441ux442ux432ux43e}

Конечное дизъюнктное объединение полуоткрытых прямоугольников:

\[
A=\bigsqcup_{j=1}^m I_j.
\]

\subsubsection{Мера Жордана на элементарных
множествах}\label{ux43cux435ux440ux430-ux436ux43eux440ux434ux430ux43dux430-ux43dux430-ux44dux43bux435ux43cux435ux43dux442ux430ux440ux43dux44bux445-ux43cux43dux43eux436ux435ux441ux442ux432ux430ux445}

Сначала на прямоугольнике:

\[
v(I)=\prod_{k=1}^d(b_k-a_k),
\]

затем на элементарном множестве:

\[
v(A)=\sum_{j=1}^m v(I_j),\qquad A=\bigsqcup_{j=1}^m I_j.
\]

\subsection{4. Основные теоремы и
свойства}\label{ux43eux441ux43dux43eux432ux43dux44bux435-ux442ux435ux43eux440ux435ux43cux44b-ux438-ux441ux432ux43eux439ux441ux442ux432ux430}

\subsubsection{Свойство 1. Алгебра замкнута относительно конечных
пересечений и
разностей}\label{ux441ux432ux43eux439ux441ux442ux432ux43e-1.-ux430ux43bux433ux435ux431ux440ux430-ux437ux430ux43cux43aux43dux443ux442ux430-ux43eux442ux43dux43eux441ux438ux442ux435ux43bux44cux43dux43e-ux43aux43eux43dux435ux447ux43dux44bux445-ux43fux435ux440ux435ux441ux435ux447ux435ux43dux438ux439-ux438-ux440ux430ux437ux43dux43eux441ux442ux435ux439}

Если \(\mathcal A\) - алгебра и \(A,B\in\mathcal A\), то

\[
A\cap B\in\mathcal A,\qquad A\setminus B\in\mathcal A.
\]

\subsubsection{Доказательство}\label{ux434ux43eux43aux430ux437ux430ux442ux435ux43bux44cux441ux442ux432ux43e}

По законам де Моргана

\[
A\cap B=\overline{\overline{A}\cup \overline{B}}.
\]

Так как \(\overline{A},\overline{B}\in\mathcal A\), то
\(\overline{A}\cup \overline{B}\in\mathcal A\), а значит и его
дополнение лежит в \(\mathcal A\). Далее

\[
A\setminus B=A\cap \overline{B},
\]

а пересечение двух множеств из \(\mathcal A\) уже доказано. \(\square\)

\subsubsection{\texorpdfstring{Свойство 2. \(\sigma\)-алгебра замкнута
относительно счетных
пересечений}{Свойство 2. \textbackslash sigma-алгебра замкнута относительно счетных пересечений}}\label{ux441ux432ux43eux439ux441ux442ux432ux43e-2.-sigma-ux430ux43bux433ux435ux431ux440ux430-ux437ux430ux43cux43aux43dux443ux442ux430-ux43eux442ux43dux43eux441ux438ux442ux435ux43bux44cux43dux43e-ux441ux447ux435ux442ux43dux44bux445-ux43fux435ux440ux435ux441ux435ux447ux435ux43dux438ux439}

Если \(A_n\in\mathcal F\) для всех \(n\), то

\[
\bigcap_{n=1}^{\infty}A_n\in\mathcal F.
\]

\subsubsection{Доказательство}\label{ux434ux43eux43aux430ux437ux430ux442ux435ux43bux44cux441ux442ux432ux43e-1}

Снова используем формулу де Моргана:

\[
\bigcap_{n=1}^{\infty}A_n=\overline{\left(\bigcup_{n=1}^{\infty}\overline{A_n}\right)}.
\]

Каждое \(\overline{A_n}\) лежит в \(\mathcal F\), счетное объединение
лежит в \(\mathcal F\), и дополнение тоже лежит в \(\mathcal F\).
\(\square\)

\subsubsection{Свойство 3. Конечные дизъюнктные объединения элементов
полуалгебры образуют
алгебру}\label{ux441ux432ux43eux439ux441ux442ux432ux43e-3.-ux43aux43eux43dux435ux447ux43dux44bux435-ux434ux438ux437ux44aux44eux43dux43aux442ux43dux44bux435-ux43eux431ux44aux435ux434ux438ux43dux435ux43dux438ux44f-ux44dux43bux435ux43cux435ux43dux442ux43eux432-ux43fux43eux43bux443ux430ux43bux433ux435ux431ux440ux44b-ux43eux431ux440ux430ux437ux443ux44eux442-ux430ux43bux433ux435ux431ux440ux443}

Пусть \(\mathcal S\) - полуалгебра и \(\Omega\in\mathcal S\). Рассмотрим
класс

\[
\mathcal A(\mathcal S)=\left\{\bigsqcup_{k=1}^n S_k:\ S_k\in\mathcal S,\ n<\infty\right\}.
\]

Тогда \(\mathcal A(\mathcal S)\) является алгеброй.

\subsubsection{Доказательство}\label{ux434ux43eux43aux430ux437ux430ux442ux435ux43bux44cux441ux442ux432ux43e-2}

Во-первых, \(\Omega\in\mathcal A(\mathcal S)\), потому что
\(\Omega\in\mathcal S\). Во-вторых, конечное объединение двух множеств
из \(\mathcal A(\mathcal S)\) можно сделать дизъюнктным разбиением,
последовательно вычитая уже учтенные части. Полуалгебра гарантирует, что
каждая разность раскладывается в конечное дизъюнктное объединение
элементов \(\mathcal S\). Значит объединение снова лежит в
\(\mathcal A(\mathcal S)\).

Остается дополнение. Если

\[
A=\bigsqcup_{k=1}^n S_k,
\]

то \(\Omega\setminus A\) получается последовательным вычитанием
\(S_1,\ldots,S_n\) из \(\Omega\). На каждом шаге разность элемента
полуалгебры с элементом полуалгебры раскладывается в конечное
дизъюнктное объединение элементов полуалгебры. После конечного числа
шагов получаем конечное дизъюнктное объединение элементов
\(\mathcal S\). Значит \(\overline{A}\in\mathcal A(\mathcal S)\).
\(\square\)

\subsubsection{Свойство 4. Мера Дирака конечно- и
счетно-аддитивна}\label{ux441ux432ux43eux439ux441ux442ux432ux43e-4.-ux43cux435ux440ux430-ux434ux438ux440ux430ux43aux430-ux43aux43eux43dux435ux447ux43dux43e--ux438-ux441ux447ux435ux442ux43dux43e-ux430ux434ux434ux438ux442ux438ux432ux43dux430}

Если \(A_1,A_2,\ldots\) попарно не пересекаются, то

\[
\delta_x\left(\bigsqcup_{n=1}^{\infty}A_n\right)=\sum_{n=1}^{\infty}\delta_x(A_n).
\]

\subsubsection{Доказательство}\label{ux434ux43eux43aux430ux437ux430ux442ux435ux43bux44cux441ux442ux432ux43e-3}

Если \(x\notin\bigcup_n A_n\), то обе части равны \(0\). Если
\(x\in\bigcup_n A_n\), то из попарной непересекаемости существует ровно
один номер \(n_0\), для которого \(x\in A_{n_0}\). Тогда левая часть
равна \(1\), а в правой сумме ровно одно слагаемое равно \(1\),
остальные равны \(0\). Значит равенство выполнено. \(\square\)

\subsubsection{Свойство 5. Полуоткрытые прямоугольники образуют
полуалгебру}\label{ux441ux432ux43eux439ux441ux442ux432ux43e-5.-ux43fux43eux43bux443ux43eux442ux43aux440ux44bux442ux44bux435-ux43fux440ux44fux43cux43eux443ux433ux43eux43bux44cux43dux438ux43aux438-ux43eux431ux440ux430ux437ux443ux44eux442-ux43fux43eux43bux443ux430ux43bux433ux435ux431ux440ux443}

В \(\mathbb R^d\) класс

\[
\mathcal S_d=\left\{\prod_{k=1}^d[a_k,b_k):\ a_k\le b_k\right\}
\]

является полуалгеброй.

\subsubsection{Доказательство}\label{ux434ux43eux43aux430ux437ux430ux442ux435ux43bux44cux441ux442ux432ux43e-4}

Пустое множество можно получить как вырожденный прямоугольник, например
\([a,a)\times\prod_{k=2}^d[a_k,b_k)\). Пересечение двух прямоугольников

\[
I=\prod_{k=1}^d[a_k,b_k),\qquad J=\prod_{k=1}^d[c_k,d_k)
\]

равно

\[
I\cap J=\prod_{k=1}^d[\max(a_k,c_k),\min(b_k,d_k)),
\]

если все интервалы невырождены; иначе пересечение пусто. Значит
пересечение снова принадлежит \(\mathcal S_d\).

Для разности \(I\setminus J\) нужно разрезать \(I\) гиперплоскостями,
проходящими через координаты \(c_k,d_k\). В каждой координате возникает
конечное разбиение на полуинтервалы. Декартовы произведения этих
полуинтервалов дают конечное разбиение \(I\) на полуоткрытые
прямоугольники; часть из них лежит в \(J\), остальные образуют
\(I\setminus J\). Значит разность есть конечное дизъюнктное объединение
элементов \(\mathcal S_d\). \(\square\)

\subsubsection{Свойство 6. Объем Жордана корректно определен на
элементарных
множествах}\label{ux441ux432ux43eux439ux441ux442ux432ux43e-6.-ux43eux431ux44aux435ux43c-ux436ux43eux440ux434ux430ux43dux430-ux43aux43eux440ux440ux435ux43aux442ux43dux43e-ux43eux43fux440ux435ux434ux435ux43bux435ux43d-ux43dux430-ux44dux43bux435ux43cux435ux43dux442ux430ux440ux43dux44bux445-ux43cux43dux43eux436ux435ux441ux442ux432ux430ux445}

Если элементарное множество \(A\) имеет два дизъюнктных разбиения на
полуоткрытые прямоугольники,

\[
A=\bigsqcup_{i=1}^m I_i=\bigsqcup_{j=1}^n J_j,
\]

то

\[
\sum_{i=1}^m v(I_i)=\sum_{j=1}^n v(J_j).
\]

\subsubsection{Доказательство}\label{ux434ux43eux43aux430ux437ux430ux442ux435ux43bux44cux441ux442ux432ux43e-5}

Рассмотрим все непустые пересечения \(I_i\cap J_j\). Они являются
полуоткрытыми прямоугольниками. Для каждого фиксированного \(i\)
множество \(I_i\) разбивается в дизъюнктное объединение прямоугольников
\(I_i\cap J_j\), поэтому

\[
v(I_i)=\sum_j v(I_i\cap J_j).
\]

Для каждого фиксированного \(j\) аналогично

\[
v(J_j)=\sum_i v(I_i\cap J_j).
\]

Суммируя по \(i\) и \(j\), получаем

\[
\sum_i v(I_i)=\sum_{i,j}v(I_i\cap J_j)=\sum_j v(J_j).
\]

Здесь используется элементарный факт: объем прямоугольника равен сумме
объемов прямоугольников конечного разбиения. Он доказывается
перемножением длин координатных разбиений. \(\square\)

\subsubsection{Свойство 7. Мера Жордана конечно-аддитивна на
элементарных
множествах}\label{ux441ux432ux43eux439ux441ux442ux432ux43e-7.-ux43cux435ux440ux430-ux436ux43eux440ux434ux430ux43dux430-ux43aux43eux43dux435ux447ux43dux43e-ux430ux434ux434ux438ux442ux438ux432ux43dux430-ux43dux430-ux44dux43bux435ux43cux435ux43dux442ux430ux440ux43dux44bux445-ux43cux43dux43eux436ux435ux441ux442ux432ux430ux445}

Если \(A,B\) - непересекающиеся элементарные множества, то

\[
v(A\sqcup B)=v(A)+v(B).
\]

\subsubsection{Доказательство}\label{ux434ux43eux43aux430ux437ux430ux442ux435ux43bux44cux441ux442ux432ux43e-6}

Представим \(A\) и \(B\) как конечные дизъюнктные объединения
полуоткрытых прямоугольников. Тогда \(A\sqcup B\) уже представлено как
конечное дизъюнктное объединение всех этих прямоугольников. По
определению меры Жордана на элементарных множествах

\[
v(A\sqcup B)=\sum_i v(I_i)+\sum_j v(J_j)=v(A)+v(B).
\]

Для конечного числа попарно непересекающихся элементарных множеств
доказательство такое же. \(\square\)

\subsection{5. Примеры}\label{ux43fux440ux438ux43cux435ux440ux44b}

\subsubsection{\texorpdfstring{Пример 1. Минимальная и максимальная
\(\sigma\)-алгебры}{Пример 1. Минимальная и максимальная \textbackslash sigma-алгебры}}\label{ux43fux440ux438ux43cux435ux440-1.-ux43cux438ux43dux438ux43cux430ux43bux44cux43dux430ux44f-ux438-ux43cux430ux43aux441ux438ux43cux430ux43bux44cux43dux430ux44f-sigma-ux430ux43bux433ux435ux431ux440ux44b}

На любом \(\Omega\) класс \(\{\varnothing,\Omega\}\) - минимальная
\(\sigma\)-алгебра. Класс \(2^\Omega\) - максимальная
\(\sigma\)-алгебра.

\subsubsection{\texorpdfstring{Пример 2. Алгебра, которая не является
\(\sigma\)-алгеброй}{Пример 2. Алгебра, которая не является \textbackslash sigma-алгеброй}}\label{ux43fux440ux438ux43cux435ux440-2.-ux430ux43bux433ux435ux431ux440ux430-ux43aux43eux442ux43eux440ux430ux44f-ux43dux435-ux44fux432ux43bux44fux435ux442ux441ux44f-sigma-ux430ux43bux433ux435ux431ux440ux43eux439}

Пусть \(\Omega=\mathbb N\). Рассмотрим класс

\[
\mathcal A=\{A\subset\mathbb N:\ A \text{ конечно или } \overline{A} \text{ конечно}\}.
\]

Это алгебра: дополнение конечного множества ко-конечно, дополнение
ко-конечного множества конечно; конечные объединения конечных множеств
конечны, а объединение с ко-конечным множеством ко-конечно.

Но \(\mathcal A\) не является \(\sigma\)-алгеброй. Одноточечные
множества \(\{2\},\{4\},\{6\},\ldots\) конечны, значит лежат в
\(\mathcal A\). Их счетное объединение - множество четных чисел. Оно
бесконечно, и его дополнение тоже бесконечно. Значит множество четных
чисел не лежит в \(\mathcal A\).

\subsubsection{Пример 3. Полуалгебра полуинтервалов на
прямой}\label{ux43fux440ux438ux43cux435ux440-3.-ux43fux43eux43bux443ux430ux43bux433ux435ux431ux440ux430-ux43fux43eux43bux443ux438ux43dux442ux435ux440ux432ux430ux43bux43eux432-ux43dux430-ux43fux440ux44fux43cux43eux439}

Класс

\[
\mathcal S=\{[a,b):\ -\infty<a\le b<\infty\}
\]

является полуалгеброй на \(\mathbb R\) в локальном смысле. Пересечение
двух полуинтервалов снова полуинтервал:

\[
[a,b)\cap[c,d)=[\max(a,c),\min(b,d)),
\]

если правый конец больше левого; иначе это пустое множество. Разность
\([a,b)\setminus[c,d)\) является объединением не более чем двух
полуинтервалов:

\[
[a,b)\setminus[c,d)= [a,\min(b,c))\sqcup[\max(a,d),b),
\]

после удаления пустых частей.

\subsubsection{Пример 4. Мера
Дирака}\label{ux43fux440ux438ux43cux435ux440-4.-ux43cux435ux440ux430-ux434ux438ux440ux430ux43aux430}

Пусть \(\Omega=\mathbb R\), \(x=0\). Тогда

\[
\delta_0((-1,1))=1,\qquad \delta_0((1,2))=0.
\]

Если случайная величина почти наверное равна \(0\), ее распределение
есть \(\delta_0\).

\subsubsection{Пример 5. Мера Жордана на прямой и
плоскости}\label{ux43fux440ux438ux43cux435ux440-5.-ux43cux435ux440ux430-ux436ux43eux440ux434ux430ux43dux430-ux43dux430-ux43fux440ux44fux43cux43eux439-ux438-ux43fux43bux43eux441ux43aux43eux441ux442ux438}

На прямой

\[
v([a,b))=b-a.
\]

Если

\[
A=[0,1)\sqcup[2,5),
\]

то

\[
v(A)=1+3=4.
\]

В \(\mathbb R^2\) для прямоугольника \([a,b)\times[c,d)\):

\[
v([a,b)\times[c,d))=(b-a)(d-c).
\]

\subsection{6. Доказательства и
обоснования}\label{ux434ux43eux43aux430ux437ux430ux442ux435ux43bux44cux441ux442ux432ux430-ux438-ux43eux431ux43eux441ux43dux43eux432ux430ux43dux438ux44f}

Для устного экзамена удобно держать три доказательные линии.

\textbf{Первая линия: операции над системами множеств.}

Все свойства алгебр и \(\sigma\)-алгебр доказываются через дополнения и
формулы де Моргана:

\[
\overline{\left(\bigcup_n A_n\right)}=\bigcap_n \overline{A_n},\qquad
\overline{\left(\bigcap_n A_n\right)}=\bigcup_n \overline{A_n}.
\]

Если система замкнута относительно дополнения и объединения, она
автоматически замкнута относительно пересечения. Если объединение
разрешено счетное, то и пересечение разрешено счетное.

\textbf{Вторая линия: полуалгебра как удобный класс разбиений.}

Полуалгебра не обязана быть замкнута относительно объединения. Но она
устроена так, что при вычитании одного простого множества из другого
получается конечное число простых множеств. Поэтому из полуалгебры можно
построить алгебру конечными дизъюнктными объединениями.

\textbf{Третья линия: меры Дирака и Жордана.}

Для Дирака доказательство аддитивности сводится к вопросу: где находится
точка \(x\)? В дизъюнктном объединении она может лежать только в одном
слагаемом.

Для Жордана главное - корректность определения. Если одно и то же
элементарное множество разбили на прямоугольники двумя способами, нужно
перейти к общему измельчению всеми попарными пересечениями
прямоугольников. Тогда обе суммы объемов сводятся к одной и той же
двойной сумме.

\subsection{7. Что написать на
доске}\label{ux447ux442ux43e-ux43dux430ux43fux438ux441ux430ux442ux44c-ux43dux430-ux434ux43eux441ux43aux435}

\begin{enumerate}
\def\labelenumi{\arabic{enumi}.}
\tightlist
\item
  Цепочка систем:
\end{enumerate}

\[
\text{полуалгебра}\longrightarrow \text{алгебра конечных дизъюнктных объединений}\longrightarrow \sigma\text{-алгебра}.
\]

\begin{enumerate}
\def\labelenumi{\arabic{enumi}.}
\setcounter{enumi}{1}
\tightlist
\item
  Определения:
\end{enumerate}

\[
A,B\in\mathcal A\Rightarrow \overline{A},A\cup B\in\mathcal A,
\]

\[
A_n\in\mathcal F\Rightarrow \bigcup_{n=1}^{\infty}A_n\in\mathcal F.
\]

\begin{enumerate}
\def\labelenumi{\arabic{enumi}.}
\setcounter{enumi}{2}
\tightlist
\item
  Полуалгебра:
\end{enumerate}

\[
A,B\in\mathcal S\Rightarrow A\cap B\in\mathcal S,\qquad
A\setminus B=\bigsqcup_{k=1}^n C_k,\ C_k\in\mathcal S.
\]

\begin{enumerate}
\def\labelenumi{\arabic{enumi}.}
\setcounter{enumi}{3}
\tightlist
\item
  Дирак:
\end{enumerate}

\[
\delta_x(A)=\mathbf 1_A(x).
\]

\begin{enumerate}
\def\labelenumi{\arabic{enumi}.}
\setcounter{enumi}{4}
\tightlist
\item
  Жордан:
\end{enumerate}

\[
v\left(\prod_{k=1}^d[a_k,b_k)\right)=\prod_{k=1}^d(b_k-a_k),\qquad
v\left(\bigsqcup_{j=1}^m I_j\right)=\sum_{j=1}^m v(I_j).
\]

\subsection{8. Типичные дополнительные
вопросы}\label{ux442ux438ux43fux438ux447ux43dux44bux435-ux434ux43eux43fux43eux43bux43dux438ux442ux435ux43bux44cux43dux44bux435-ux432ux43eux43fux440ux43eux441ux44b}

\textbf{Вопрос. Почему в определении алгебры достаточно требовать
замкнутость относительно дополнения и объединения, но не пересечения?}

Потому что пересечение выражается через дополнение и объединение:

\[
A\cap B=\overline{\overline{A}\cup \overline{B}}.
\]

\textbf{Вопрос. Чем алгебра отличается от \(\sigma\)-алгебры?}

Алгебра замкнута относительно конечных объединений, а \(\sigma\)-алгебра
- относительно счетных объединений. Например, конечные и ко-конечные
подмножества \(\mathbb N\) образуют алгебру, но не \(\sigma\)-алгебру.

\textbf{Вопрос. Почему полуинтервалы \([a,b)\) берут именно
полуоткрытыми?}

Потому что они хорошо дизъюнктно склеиваются без пересечения границ:

\[
[a,c)=[a,b)\sqcup[b,c).
\]

Это удобно для задания длины и дальнейшего продолжения меры.

\textbf{Вопрос. Является ли полуалгебра алгеброй?}

Не обязательно. Полуинтервалы \([a,b)\) не замкнуты относительно
объединения: \([0,1)\cup[2,3)\) не является одним полуинтервалом. Но
конечные дизъюнктные объединения таких полуинтервалов уже образуют
алгебру элементарных множеств.

\textbf{Вопрос. Почему мера Дирака является вероятностной мерой?}

Потому что \(\delta_x(\Omega)=1\), \(\delta_x(A)\ge0\), и она
счетно-аддитивна на дизъюнктных объединениях.

\textbf{Вопрос. Почему мера Жордана не решает всю теорию меры?}

Потому что она хорошо работает для элементарных и жорданово измеримых
множеств, но не образует достаточно богатой \(\sigma\)-аддитивной теории
на всех множествах, которые возникают как пределы. Для этого строят меру
Лебега.

\subsection{9. Типичные
ошибки}\label{ux442ux438ux43fux438ux447ux43dux44bux435-ux43eux448ux438ux431ux43aux438}

\begin{itemize}
\tightlist
\item
  Говорить, что полуалгебра замкнута относительно объединений. Это
  неверно в общем случае.
\item
  Забывать, что алгебра требует только конечные операции, а
  \(\sigma\)-алгебра - счетные.
\item
  Путать \(\varnothing\in\mathcal S\) и \(\Omega\in\mathcal S\): для
  полуалгебры обязательно нужно пустое множество, а наличие \(\Omega\)
  зависит от принятого варианта определения.
\item
  Считать, что мера Жордана задана на всех подмножествах
  \(\mathbb R^d\). Она изначально задана на элементарных множествах, а
  затем на жорданово измеримых множествах.
\item
  Не проверять корректность определения \(v(A)=\sum_jv(I_j)\): одно и то
  же элементарное множество может иметь разные разбиения.
\item
  Называть любое неотрицательное отображение мерой. Нужна хотя бы
  конечная аддитивность, а в вероятностной теории дальше потребуется
  счетная аддитивность.
\end{itemize}

\subsection{10. Связи с другими
билетами}\label{ux441ux432ux44fux437ux438-ux441-ux434ux440ux443ux433ux438ux43cux438-ux431ux438ux43bux435ux442ux430ux43cux438}

\begin{itemize}
\tightlist
\item
  Билет 02: счетная аддитивность и вероятностные пространства используют
  \(\sigma\)-алгебру как систему событий.
\item
  Билет 03: продолжение меры начинается с меры на полуалгебре и строит
  меру на алгебре и \(\sigma\)-алгебре.
\item
  Билет 04: борелевская \(\sigma\)-алгебра - главный пример
  \(\sigma\)-алгебры в метрических пространствах.
\item
  Билет 05: измеримость случайной величины означает, что прообразы
  борелевских множеств лежат в \(\mathcal F\).
\item
  Билет 06: распределение случайной величины - это мера на пространстве
  значений.
\end{itemize}

\subsection{11. Минимум на
удовлетворительно}\label{ux43cux438ux43dux438ux43cux443ux43c-ux43dux430-ux443ux434ux43eux432ux43bux435ux442ux432ux43eux440ux438ux442ux435ux43bux44cux43dux43e}

Нужно уметь строго сформулировать:

\begin{enumerate}
\def\labelenumi{\arabic{enumi}.}
\tightlist
\item
  определение алгебры;
\item
  определение \(\sigma\)-алгебры;
\item
  определение полуалгебры;
\item
  формулу меры Дирака;
\item
  формулу объема прямоугольника для меры Жордана.
\end{enumerate}

Минимальный ответ:

\begin{quote}
Алгебра замкнута относительно конечных операций над множествами,
\(\sigma\)-алгебра - относительно счетных объединений, полуалгебра
замкнута относительно пересечений, а разности в ней раскладываются в
конечные дизъюнктные объединения. Мера Дирака задается формулой
\(\delta_x(A)=\mathbf 1_A(x)\), мера Жордана на прямоугольнике -
формулой \(v(\prod[a_k,b_k))=\prod(b_k-a_k)\).
\end{quote}

\subsection{12. Минимум на
хорошо}\label{ux43cux438ux43dux438ux43cux443ux43c-ux43dux430-ux445ux43eux440ux43eux448ux43e}

Кроме минимума нужно доказать:

\begin{enumerate}
\def\labelenumi{\arabic{enumi}.}
\tightlist
\item
  что алгебра замкнута относительно пересечений и разностей;
\item
  что \(\sigma\)-алгебра замкнута относительно счетных пересечений;
\item
  что полуинтервалы \([a,b)\) образуют полуалгебру;
\item
  что мера Дирака аддитивна.
\end{enumerate}

Также нужно привести пример алгебры, не являющейся \(\sigma\)-алгеброй:
конечные и ко-конечные подмножества \(\mathbb N\).

\subsection{13. Что добавить для
отлично}\label{ux447ux442ux43e-ux434ux43eux431ux430ux432ux438ux442ux44c-ux434ux43bux44f-ux43eux442ux43bux438ux447ux43dux43e}

Для отличного ответа нужно объяснить роль полуалгебры в построении меры:

\[
\mathcal S\longrightarrow \mathcal A(\mathcal S)\longrightarrow \sigma(\mathcal S).
\]

Также стоит доказать корректность меры Жордана через общее измельчение
двух разбиений. Это показывает, что формула

\[
v(A)=\sum_j v(I_j)
\]

не зависит от выбранного представления элементарного множества \(A\).

Полезно добавить предупреждение: мера Жордана - предшественник меры
Лебега, но она недостаточно устойчива для предельных операций, поэтому в
вероятности и анализе переходят к \(\sigma\)-аддитивной мере Лебега.

\subsection{14. Устная версия ответа на 5
минут}\label{ux443ux441ux442ux43dux430ux44f-ux432ux435ux440ux441ux438ux44f-ux43eux442ux432ux435ux442ux430-ux43dux430-5-ux43cux438ux43dux443ux442}

Сначала я фиксирую пространство \(\Omega\) и говорю, что события - это
не обязательно все подмножества \(\Omega\), а элементы некоторой системы
множеств. Для конечных операций достаточно алгебры:
\(\Omega\in\mathcal A\), вместе с \(A\) лежит \(\overline{A}\), вместе с
\(A,B\) лежит \(A\cup B\). Через де Моргана сразу получаю
\(A\cap B\in\mathcal A\) и \(A\setminus B\in\mathcal A\).

Дальше усиливаю определение: \(\sigma\)-алгебра - это алгебра, замкнутая
относительно счетных объединений. Тогда она замкнута и относительно
счетных пересечений, потому что

\[
\bigcap_n A_n=\overline{\left(\bigcup_n \overline{A_n}\right)}.
\]

После этого объясняю полуалгебру. Это система простых множеств, где
пересечение остается простым, а разность раскладывается в конечное
дизъюнктное объединение простых множеств. Главный пример - полуинтервалы
\([a,b)\) на прямой и полуоткрытые прямоугольники в \(\mathbb R^d\).

Затем перехожу к мерам. Мера Дирака в точке \(x\) задается формулой
\(\delta_x(A)=\mathbf 1_A(x)\). Она аддитивна, потому что в дизъюнктном
объединении точка \(x\) может лежать максимум в одном слагаемом.

Мера Жордана - это длина, площадь или объем элементарных множеств. На
прямоугольнике \(I=\prod[a_k,b_k)\) она равна \(v(I)=\prod(b_k-a_k)\).
На конечном дизъюнктном объединении прямоугольников она равна сумме
объемов. Корректность доказывается общим измельчением двух разбиений.

В конце связываю билет со следующими: в билете 2 нужна счетная
аддитивность, в билете 3 - продолжение меры с полуалгебры, в билете 4 -
борелевские \(\sigma\)-алгебры.

\subsection{15. Если осталось 2 минуты перед
экзаменом}\label{ux435ux441ux43bux438-ux43eux441ux442ux430ux43bux43eux441ux44c-2-ux43cux438ux43dux443ux442ux44b-ux43fux435ux440ux435ux434-ux44dux43aux437ux430ux43cux435ux43dux43eux43c}

Запомнить пять формул:

\[
A\cap B=\overline{\overline{A}\cup \overline{B}},
\]

\[
\bigcap_n A_n=\overline{\left(\bigcup_n \overline{A_n}\right)},
\]

\[
A\setminus B=\bigsqcup_{k=1}^n C_k\quad\text{для полуалгебры},
\]

\[
\delta_x(A)=\mathbf 1_A(x),
\]

\[
v\left(\prod_{k=1}^d[a_k,b_k)\right)=\prod_{k=1}^d(b_k-a_k).
\]

Главная ловушка: алгебра - конечные операции, \(\sigma\)-алгебра -
счетные операции, полуалгебра - не алгебра и не обязана быть замкнута по
объединениям.

\newpage

\end{document}
